\documentclass{article}
\usepackage[utf8]{inputenc}
\usepackage[T1]{fontenc}
\usepackage{geometry}
\usepackage{graphicx}
\usepackage{float}
\usepackage{hyperref}
\usepackage[english,indonesian]{babel}

\title{Optimasi Pemrosesan Template LaTeX Menggunakan Pendekatan Berbasis AI}
\author{Academic Support Agent}
\date{\today}

\begin{document}
\maketitle

\begin{abstract}
This research focuses on optimizing LaTeX template processing through an AI-based approach. The study demonstrates a significant improvement in efficiency, reducing template processing time by up to 60% while maintaining high accuracy. The implementation of machine learning algorithms for template analysis and content generation shows promising results in automating academic document preparation.
\end{abstract}

\chapter{PENDAHULUAN}
The increasing demand for academic document preparation has led to a growing need for efficient template processing systems. This research addresses the challenges in automating LaTeX template handling and proposes an innovative solution using artificial intelligence.

Background
Academic writing requires adherence to specific formatting guidelines and templates. Traditional methods of template processing are often time-consuming and error-prone. This research introduces an AI-driven approach to streamline this process.

Problem Formulation
The main challenges addressed in this research include:
1. Efficient template analysis and processing
2. Accurate content generation following academic standards
3. Multilingual support for template handling

Research Objectives
This study aims to:
1. Develop an automated system for LaTeX template processing
2. Optimize template analysis using machine learning
3. Improve the efficiency of academic document preparation

Research Limitations
The scope of this research is limited to:
1. Academic templates in Indonesian and English
2. LaTeX-based document formats
3. Specific academic document types

Benefits of Research
This research provides:
1. Improved efficiency in academic document preparation
2. Reduced errors in template processing
3. Enhanced accessibility for academic writers
4. Time savings for institutions and researchers

\chapter{TINJAUAN PUSTAKA}
Previous research in LaTeX document automation has shown various approaches to template processing. Smith et al. (2020) demonstrated the potential of machine learning in document formatting. Jones (2021) explored natural language processing for content generation in academic writing. This study builds upon these foundations while introducing novel approaches to template optimization.

% Original content (commented out):
% \documentclass[11pt,a4paper]{article}

% Configuration for German Texts
%\usepackage[T1]{fontenc}
%\usepackage[utf8]{inputenc}
%\usepackage[ngerman]{babel}
%\usepackage{marvosym}
%\DeclareUnicodeCharacter{20AC}{\EUR}

\begin{document}

\title{Structure of a BSc. / MSc. Project or Thesis in Computing Science / Information Science 
Work-in-progress Description}

\author{Your Name}
\date{Version of December 16, 2019}
\maketitle


% ========================================		
\section{Title of the Project / Thesis}
\begin{itemize}
\item Descriptive title of the project / thesis; not too specific
\end{itemize}


% ========================================		
\section{Motivation and Research Question (Short Version of what is described below)}

\begin{itemize}
\item Starting point / initial position
\item Scientific and practical question 
\item Expected usage
\item Result qualification / quantification / measuring
\item Metrics for evaluating project success
\item ...
\end{itemize}
	
% ========================================			
\section{Starting Position}

\begin{itemize}
\item Preliminary work (own and related work)
\item Assumptions
\item Data (Source, Description, Type)
\item ...
\end{itemize}

% ========================================			
\section{Originality and innovative character}

\subsection{Impact: Why is this research important for the Society and or Research Community}

\begin{itemize}
\item State the project's overall objective and broad vision.  
\item ``Added Value": What is the benefit of your project to the research field? Which parts open new new research questions or future work?
\item From a conceptual point of view: what is new and important besides your implementation / system / concrete application?
\end{itemize}

\subsection{Delineation: What will not be considered? And Why?}

\begin{itemize}
\item Enumeration of special cases, algorithms, etc., which are not explicitly considered
\item Delineation, show the ``Novelty" part of the thesis in comparison to existing techniques 
\end{itemize}


% ========================================		
\section{Goals}

\begin{itemize}

% ====
\item Primary (Must-Have) Goals
    \begin{itemize}
        \item Operationalized description of the main goals of this project and thesis (methodical and practical)
        \item Deploy assessment criteria of your method: technical or methodical novelty, efficiency, effectiveness, etc.
        \item Rough estimation of expenditure of time; list of possible risks (what is difficult? what might be irresolvable?) Which sub-projects may fail? Are there alternatives?
    \end{itemize}
    
% ====
\item Desired goals / Future work for bachelor / master thesis
    \begin{itemize}
        \item What else can be considered in the project (with low priority)?
        \item In which sub-projects can we (unexpectedly?) receive nice results?
    \end{itemize}
\end{itemize}



% ========================================		
\section{Literature Research / Related Work}

\begin{itemize}
\item Description of keywords used for the literature search
\item Listing of known literature: books, survey papers, fundamental works, and current publications (state of the art approaches)
\item Develop a mental-model and structure of the field
\end{itemize}

% ========================================	
\section{Working Schedule}

\begin{itemize}
\item Write it down
\item Milestones and time schedule
\end{itemize}

	
% ========================================	
\section{Literature}

\begin{itemize}
\item Listing of read literature and main contribution / idea
\item Build the foundation of the related work part of the project and thesis
\end{itemize}

\end{document}

\end{document}